\chapter{Introduction}
\label{ch:introduction}

\section{Motivation}

Many decision problems in security, negotiation, finance, and multi-agent
interaction share three properties that standard fully observed reinforcement
learning benchmarks underrepresent.

\textbf{Partial observability.} Important state variables remain hidden until
late in an interaction. An agent must infer latent structure from noisy clues
rather than read full state directly.

\textbf{Simultaneous action selection.} Both sides commit to actions before the
environment resolves the step. This differs from chess-like turn-taking and
creates a distinct learning problem around anticipating opponent intent under
uncertainty.

\textbf{Combinatorial structure.} The space of valid configurations grows
quickly. A model that memorizes frequent patterns may appear strong on average
yet fail on rare but valid combinations held out from training.

Transformers are a natural fit for these settings because they aggregate evidence
across long interaction histories with causal attention. Offline learning is
appropriate when high-quality human trajectories are available at scale and the
deployment goal is a single forward pass without lookahead search or simulator
queries at inference time.

\section{Problem Statement}

Given a sequence of first-person observations from past turns, we seek a
model that
\begin{enumerate}[label=(\roman*)]
  \item maintains a useful belief over hidden opponent attributes,
  \item uses that belief to choose strong actions in new interactions, and
  \item generalizes to opponent configurations outside the training distribution.
\end{enumerate}

Formally, each turn produces an observation $o_t$ derived from public state.
Hidden variables $z_t$ include unrevealed moves, items, abilities, and
other private attributes. Actions $a_t$ are chosen under simultaneous-move
constraints. The model receives offline trajectories
$\{(o_{1:T}, a_{1:T}, z_{1:T})\}$ reconstructed from full logs where $z$ is
known post hoc but not visible to the agent at decision time.

The central thesis question is whether explicit belief supervision and
transformer-based history modeling improve both belief accuracy and policy
performance under combinatorial distribution shift.

\section{Scope and Non-Goals}

This thesis measures offline imitation and belief metrics. It does
\emph{not} claim that the resulting agent is competitive against human ladder
players. Live win rate is a separate play-strength question. Offline action
match of roughly twenty percent on hundreds of classes is far above chance and
still leaves most turns incorrect, so weak live performance against skilled
humans is expected and does not invalidate the scientific measurements below.

Pok\'emon is used only as a benchmark. The contribution is an empirical study of
belief inference and combinatorial generalization under offline learning, not a
product bot or a search-based competitive agent.

\section{Research Questions}

\begin{description}
  \item[RQ1 (Belief inference)] Can a causal transformer predict hidden opponent
    attributes more accurately than classical and neural baselines?
  \item[RQ2 (Belief usefulness)] Does improved belief accuracy translate into
    better action selection under behavioral cloning and offline policy metrics?
  \item[RQ3 (Combinatorial generalization)] Under held-out team splits,
    do belief-augmented policies degrade more gracefully than action-only models?
  \item[RQ4 (Architecture)] How much gain comes from within-turn structure,
    cross-turn causal attention, and auxiliary belief heads versus simpler
    baselines?
\end{description}

\section{Contributions}

This thesis makes four contributions.
\begin{enumerate}
  \item A structured tokenization and belief supervision framework for
    partially observed simultaneous-action games.
  \item Empirical evidence that transformers and classical baselines learn
    nontrivial beliefs over hidden opponent moves from offline human play, and
    that a latest-turn MLP can outperform a history transformer on that target.
  \item A combinatorial generalization protocol with held-out opponent team keys,
    showing that team-level OOD gaps shrink primarily with data scale rather than
    with explicit matchup or belief modules.
  \item An open-source implementation built on Metamon, poke-env, and Gymnasium
    compatible tooling, including export, training, probing, and OOD evaluation
    scripts.
\end{enumerate}

\section{Evaluation Domain}

We evaluate on competitive Pok\'emon battles reconstructed from human spectator
replays via the Metamon dataset \citep{grigsby2023metamon,grigsby2025metamon}.
Pok\'emon combines partial observability, simultaneous move commitment, and
combinatorial team structure in one domain with large-scale offline data. Gen~9
OU team preview reveals all opponent species at battle start, so species belief
is empty by design and active-opponent move, item, ability, and tera labels are
the primary belief targets. The methods apply to any domain with similar
structure.

\section{Thesis Organization}

\Cref{ch:background} reviews POMDPs, transformers in decision making, and
offline RL. \Cref{ch:benchmark} briefly describes the benchmark environment.
\Cref{ch:representation,ch:model} present tokenization and the belief
transformer. \Cref{ch:experiments,ch:results,ch:analysis} cover the experiment
plan, results, and interpretability analysis. \Cref{ch:limitations,ch:conclusion}
discuss limitations and future work.
